\scanpage{61}
\mspage{10.}

et $\psi\varphi\sim\mathrm{id}_C$ (on trouve même une flèche $\psi\varphi\unsure{\text{---}|\longrightarrow}\mathrm{id}_C$), Ainsi $\varphi$ est un homotopisme, donc une $W$-équivalence (cor. 3 p. 3), donc $C$ est $W$-asphérique, $C'$ l'étant.

La propreté de $p_S$ équivaut à la lissité de $p_{(S^\circ)}$, elle est donc prouvée du même coup.

Ainsi, le lemme de la page 7 est prouvé, moyennant l'axiome \textcircled{4} qui assure qu'une flèche qui est à la fois $W$-lisse et $W$-propre, satisfait à la condition très forte \inserted{a)} des $W$-fibrations.

\margin{\uline{NB}}

En fait, il suffirait, au lieu de l'axiome \textcircled{4} sous \deleted{\unsure{sa}} forme forte, de supposer qu'une $f:X\longrightarrow Y$ \deleted{qui} est à la fois ``homotopiquement lisse'' et ``homotopiquement propre'', i.e. les fibres des $X_{/x}\longrightarrow Y_{/y}$ et des $ {}_{x\backslash}X\longrightarrow{}_{y\backslash}Y$ sont \unsure{contractiles|contractibles} (donc $W$-asphériques pour tout $W$ satisfaisant \textcircled{1}\textcircled{2}\textcircled{3}), soit $W$-fibrante forte. Énonçons quand même ce que nous avons prouvé.

% slanted
\margin{pour \textcircled{2}, il \unsure{le} suffit sous la forme $W(2\,\mathrm{bis})$}

\uline{Proposition 4} Si $W$ satisfait \textcircled{1}\textcircled{2}\textcircled{3} (nous \unsure{les supposons|le supposons} tjrs par la suite) plus \textcircled{4} (sous une forme un peu affaiblie suffirait, cf remarque précédente), on a ceci : Soit $f:X\longrightarrow Y$ une flèche dans $\mathrm{Cat}_{/S}$, avec $X,Y$ des $W$-fibrés sur $S$. \deleted{Alors} Si $f\in W$, alors $f_s:X_s\longrightarrow Y_s$ est dans $W$ pour tout $s\in S$.

Je me demande si j'en sais assez pour pouvoir prouver
\[
W\subset W_\infty.
\]

Il faut quand même éliminer \unsure{du|le} cas \unsure{canularesque|cauchemardesque}, \unsure{horreur|hmmm} $W=\mathrm{Fl}(\mathrm{Cat})$, donc mettons l'axiome

% slanted
\margin{$W(5)$ Axiome de connexité. Voir à ce sujet les cor. page 22}

\textcircled{5} Si $f:X\longrightarrow Y$ est dans $W$, $\pi_0(f):\pi_0(X)\longrightarrow\pi_0(Y)$ bijectif.

\deleted{Essayons de montrer que si $f:X\longrightarrow Y$ est dans $W$, il est dans $W_\infty$. Soit $I=\pi_0(X)\simeq\pi_0(Y)$, considérons}

\deleted{\diagram{\begin{tikzcd}[ampersand replacement=\&,row sep=tiny,column sep=small]
X\arrow[rr,"f"]\arrow[dr] \& \& Y\arrow[dl]\\
\& I
\end{tikzcd}}}

% archive sheet number: 61
